\section{Proprioception and Path Adaptation}
\label{sec:proprioception}

The open-loop control strategy and kinematic chain validated in Section 2.2 demonstrated the system's theoretical
correctness under ideal conditions. However, this approach relies on the strict assumption of a perfectly flat
and predictable terrain ($z_{ground} = p_{z,0}$). In a real-world unstructured environment, ground elevation
varies unpredictably. Without feedback, a purely time-driven limb would either strike the ground prematurely
— causing high mechanical stress and potential instability — or fail to reach it entirely, compromising propulsion.

To relax this assumption and enable a "Blind Walking" \cite{brooks1989robot} capability, the system is
augmented with basic proprioceptive sensing and a reactive control layer. Admittedly, while this approach does
not fully solve the locomotion problem for extreme terrain (which would require a more structured control), it
represents the fundamental first step towards robust locomotion: the ability to detect and adapt to the ground
surface upon contact.


\subsection{Contact Sensing and Architectural Shift}

To detect the ground transition, the limb is equipped with a binary contact sensor located at the tip of the
foot. This sensor acts as a logical limit switch, providing a boolean signal $S_c$ to the controller:

$$
  S_c =
  \begin{cases}
    1 & \text{if touching ground}\\
    0 & \text{otherwise}
  \end{cases}
$$

The introduction of this sensor demands a fundamental shift in the control architecture. The path planning
problem transforms from a purely time-dependent function, $p_d(t)$, to a state-dependent system driven by
discrete events.

In the previous open-loop model, the gait phase was dictated by a global clock. In this adaptive model, the
controller must monitor $S_c$ and be capable of interrupting the Swing Phase exactly at the moment of
physical impact. Consequently, the trajectory generator cannot simply map a normalized time $t \in [0,1]$ to a
position; it must be restructured to handle asynchronous events and variable phase durations.


\subsection{Finite State Trajectory Generator}

To implement this reactive logic robustly, a custom Finite State Machine (FSM, Figure \ref{fig:fsm_logic}) was designed
\cite{cruse1990mechanisms}. This generator decouples the logical phase of the leg from the global clock,
allowing for local adaptations while maintaining synchronization through a phase-advance mechanism.

\begin{figure}[htbp]
    \centering
    \begin{tikzpicture}[
        ->,                       % Freccia su tutti i percorsi
        >=Stealth,                % Stile punta freccia moderno
        node distance=3.5cm and 4.5cm, % Distanza verticale e orizzontale
        every state/.style={      % Stile globale per i cerchi degli stati
            draw=black,
            thick,
            fill=gray!10,
            minimum size=3cm,     % Dimensione fissa per cerchi uniformi
            align=center          % Permette "a capo" dentro i cerchi
        },
        thick                     % Linee spesse
    ]

    % --- Definizione degli Stati ---

    % 1. Stance 1 (Alto-Sinistra) - Stato Iniziale
    \node[state, initial left, initial text={Start}] (S1) {Posterior\\Stance\\(Stance 1)};

    % 2. Lift (Alto-Destra) - A destra di S1
    \node[state] (Lift) [right=of S1] {Lift\\(Early Swing)};

    % 3. Touch (Basso-Destra) - Sotto a Lift
    \node[state] (Touch) [below=of Lift] {Touch\\(Late Swing)};

    % 4. Stance 2 (Basso-Sinistra) - A sinistra di Touch (e sotto S1)
    \node[state] (S2) [left=of Touch] {Anterior\\Stance\\(Stance 2)};

    % --- Transizioni (Frecce) ---

    % 1. Stance 1 -> Lift (Freccia verso destra)
    \path (S1) edge node[above, align=center] {
        $p_h \le -r$
    } (Lift);

    % 2. Lift -> Touch (Freccia verso il basso)
    \path (Lift) edge node[right, align=left] {
        $t > T_{lift}$
    } (Touch);

    % 3. Touch -> Stance 2 (Freccia verso sinistra)
    \path (Touch) edge node[below, align=center] {
        $S_c = 1$ \\
        \textit{-- or --} \\
        $t > T_{swing}$
    } (S2);

    % 4. Stance 2 -> Stance 1 (Freccia verso l'alto)
    \path (S2) edge node[left, align=right] {
        $p_h = 0$
    } (S1);

    \end{tikzpicture}
    \caption{Finite State Machine for Adaptive Trajectory Generation.}
    \label{fig:fsm_logic}
\end{figure}


\subsubsection{State 1: Posterior Stance (Stance 1)}

The initial state of the cycle. Geometrically, it corresponds to the leg moving from the central rest
position ($p_h = 0$) towards the posterior limit of the workspace ($p_h = -r$).

\newpage
The horizontal velocity is constant ($\dot{p}_h = -v$). The position is calculated via a discrete-time
integrator, which is reset upon entering the state. The vertical position $z$ is held constant at the
previously detected ground level.

Upon entering this state, the system calculates a synchronization parameter, $\delta(k)$, to compensate for any
timing drift accumulated in the previous cycle $k-1$. Recalling the duty factor $\beta$, the parameter is calculated as:
\begin{equation}
  \delta(k) = \left(1 - \frac{\beta}{2}\right) - t - \delta(k-1)
\end{equation}
where $t$ is the normalized global time counter. This ensures that the leg remains synchronized with the
robot's gait rhythm, correcting for local timing variations.

The system transitions to the Lift state when the horizontal position reaches the posterior stroke limit
($p_h \le -r$).


\subsubsection{State 2: Lift Phase (Early Swing)}

This state represents the deterministic part of the swing phase, where the leg lifts off the ground to clear
potential obstacles.

The duration of this phase, $T_{lift}$, is dynamically calculated to absorb the phase error $\delta(k)$:
\begin{equation}
  T_{lift} = T_{touch} + \delta(k) = \frac{1-\beta}{2} + \delta(k)
\end{equation}
This variable duration is the mechanism that re-synchronizes the leg with the robot's global gait cycle.

A quintic polynomial moves the foot vertically from its current height to the clearance height $H_{lift}$ over
the duration $T_{lift}$, while another quintic polynomial starts repositioning the foot horizontally over the
duration $T_{swing} = T_{lift} + T_{touch}$.

The state transitions only when the timer exceeds the allocated time ($t > T_{lift}$), meaning the foot has
reached the desired clearence height.


\subsubsection{State 3: Touch Phase (Late Swing)}

This state governs the descent of the leg. Crucially, this is where proprioceptive feedback is expected to trigger.

The timer variable is adjusted to parametrize the descent polynomial $t' = t_{norm} - T_{lift}$, guiding the
foot from $H_{lift}$ back towards the nominal ground ($z = 0$) over a fixed duration $T_{touch} = (1-\beta)/2$.
The global timer is maintained to preserve continuity with the swing phase dynamics.

The system exits this state under two distinct conditions:

\begin{enumerate}
  \item \textbf{Time-out:} The time expires ($t' > T_{touch}$ or equivalently $t > T_{swing}$), implying the leg
    completed the full trajectory without contact (e.g. a hole or flat ground).
  \item \textbf{Early Contact:} The sensor triggers ($S_c = 1$). This indicates the ground is higher
    than expected (e.g. an obstacle).
\end{enumerate}

In both cases, the FSM immediately transitions to the next state.


\subsubsection{State 4: Anterior Stance (Stance 2)}

Following the touchdown, the leg enters the anterior portion of the stance phase, moving from the anterior limit
($+r$ or the early contact point) back towards the central rest position ($p_h = 0$).

The horizontal velocity returns to $-v$. The vertical position $z$ is locked at the specific height recorded at
the moment of the Touch $\to$ Anterior Stance transition. This allows the robot to walk "on top" of obstacles rather
than trying to push through them.

The cycle completes when the leg crosses the zero-point of the workspace ($p_h = 0$), triggering the transition
back to Posterior Stance. At this crossing, the cycle index $k$ increments, and the loop repeats.


\subsection{Adaptive Simulation Results}

To validate the proprioceptive controller, the Simscape environment was modified to include an uneven terrain
profile.

Figure \ref{fig:adaptive_results} provides a direct comparison between the nominal, time-based trajectory
and the reactive, adaptive trajectory generated by the FSM in the presence of uneven terrain. For visual
clarity, these plots represent the trajectory primitives defined in the $Z\psi$-plane.


\subsubsection{Vertical Adaptation}

The bottom plot (Vertical Trajectory) clearly demonstrates the effectiveness of the proprioceptive loop. When
the limb encounters a positive obstacle (represented by the gray terrain profile), the adaptive trajectory
deviates from the nominal path exactly at the moment of impact. Unlike the nominal reference, which would
attempt to drive the foot through the terrain, the adaptive system triggers an immediate transition from the Touch
Phase to the Stance Phase. The vertical position is subsequently maintained at the obstacle's height for the
duration of the propulsion stroke, effectively allowing the robot to walk on top of the irregularity rather than
pushing against it.

\begin{figure}[htbp]
    \centering
    \includegraphics[width=0.98\textwidth]{images/adaptive_comparison.pdf}
    \caption{Adaptive Gait Simulation with the New FSM Trajectory Generator.}
    \label{fig:adaptive_results}
\end{figure}


\subsubsection{Synchronization and Phase Correction}

Despite the timing disturbances introduced by early ground contacts, the top plot (Horizontal Trajectory) shows
that the system maintains robust synchronization. While the early impact inevitably causes a minor phase lead
(visible as a slight horizontal shift between the purple and green waves), the trajectory does not drift
indefinitely. Instead, when the obstacle is no longer present, the cycle realigns with the non-adaptive ideal
trajectory. This proves that the phase synchronization mechanism successfully compensates for these delays in
subsequent cycles, keeping the limb's rhythm bounded and consistent with the global gait cycle.


\subsubsection{Continuity and Stability Considerations}

There is key consideration to address regarding the kinematic continuity at the impact point: while the
generated trajectory primitives are designed to be $C^2$ continuous, the very nature of the "Blind Walk"
implies that the exact moment of touchdown is unpredictable. Consequently, when the ground is struck earlier
than planned, the vertical velocity changes abruptly from the descent speed to zero (relative to the terrain).
Therefore, at these specific transition points, strict $C^2$ continuity cannot be guaranteed.

Secondly, while effective as a proof of concept, this adaptive generator has specific limitations. For starters,
it currently compensates only for elevated obstacles (protrusions), neglecting potential depressions or holes
in the terrain. Also, unlike the nominal generator, this reactive strategy does not inherently guarantee that
the modified path remains within the Optimal Workspace Region (OWR); without explicit saturation limits, the
adaptation could potentially request kinematically unfeasible positions.

Another important consideration to keep in mind is the dynamic shortening of the foot's contact time introduces
significant stability challenges; therefore, a rigorous balance study is essential for the future implementation
of the full robot.

In summary, this strategy is designed specifically to handle minor terrain irregularities (roughness) rather
than large obstacles, representing a solid foundation for higher level obstacle avoidance planning and control
\cite{saranli2001rhex}.